\documentclass[10pt,letterpaper]{article}
\usepackage{cornell-notes}

\title{Clause 29.08.18: Conventional Syntax Operators}
\author{}
\date{}
\setDocTitle{Clause 29.08.18: Conventional Syntax Operators}
\setDocSubtitle{C++ 2024}
\setDocOwner{C++ 2024 Cornell Notes}
\setDocDate{\today}
\setCornellCollection{C++ 2024 Cornell Notes}
\setCornellUnitType{Clause}
\setCornellUnitNumber{29.08.18}
\setCornellUnitTitle{Conventional Syntax Operators}
\hypersetup{pdftitle={Clause 29.08.18: Conventional Syntax Operators},pdfsubject={Cornell notes},pdfauthor={}}

\begin{document}
\maketitle
\begin{CornellOverviewBox}{Overview}
\textbf{Classification:} Standard library - Normative\\[2pt]
\textbf{Learning objective:} Understand the rules, terminology, and semantic consequences associated with Conventional syntax operators.
\end{CornellOverviewBox}\begin{CornellSummaryBox}{Summary}
{\sffamily\bfseries\color{Accent} Summary}\par\vspace{3pt}Section 29.8.18 focuses on Conventional syntax operators. Apply the specified interface together with its mandates, effects, result conditions, exception behavior, and complexity constraints. When applying it, preserve the distinction between mandatory requirements, recommendations, permissions, and behavior deliberately left undefined, unspecified, or implementation-defined.
\end{CornellSummaryBox}

\end{document}
